\documentclass[../main.tex]{subfiles}
\begin{document}
%==========================================TIÊU ĐỀ TẬP========================================
\begin{table}[h]
    \begin{adjustwidth}{-1cm}{}
        \begin{tabular}{>{\centering\arraybackslash}p{6cm}|>{\raggedright\arraybackslash}p{12.5cm}}
            \multirow{5}{6cm}
            {\\[-40pt]
                \flaregothic\fontsize{20pt}{20pt}\selectfont \centering
                \color{eptype}HISTOIRE\color{black}\\
                \freshbot\fontsize{72pt}{72pt}\selectfont
                \color{epnum}31\\[6pt]
                \cafeteria\fontsize{20pt}{20pt}\selectfont\color{epnum}(D-31)\color{black}
                \\[18pt]
                \color{black}
            }
             &
            {\fotskip\fontsize{18pt}{18pt}\selectfont \ruby{世}{せ}\ruby{代}{だい}を\ruby{超}{こ}えるBBQ}
            \\[6pt]
             & {\cafeteria\fontsize{18pt}{18pt}\selectfont The Cross-Generation BBQ}                                                                                                                                                                                     \\[4pt]
             & {\vnmsans\fontsize{18pt}{18pt}\selectfont Tiệc nướng xuyên thế hệ}                                                                                                                                                                                        \\[8pt]
             & {\caviar{\fontsize{13pt}{13pt}\selectfont Nhân vật: \emp{kanata} \emp{watame} \emp{towa} \emp{luna} \emp{lamy} \emp{nene} \emp{botan} \emp{polka} \emp{kson} \emp{delu} \emp{kuon2} \emp{kaiho2} \emp{ryuuhou2} \emp{suzuno2} \emp{naomi2} \emp{koneko}}}
            \\[8pt]
        \end{tabular}
    \end{adjustwidth}
\end{table}
%=========================================TRUYỆN CHÍNH========================================
\color{black}

\txtbx[2]{
    {\color{vol} \uline{Tóm tắt:}} Một chiều thứ Bảy giữa tháng Bảy đầy nắng, sân thượng tầng bảy của ký túc xá Hikawa trở thành điểm hẹn lịch sử khi các thành viên của hai thế hệ Gen 4 và Gen 5 lần đầu tiên cùng tề tựu đông đủ. Dưới sự điều phối của anh Kuon cùng sự hỗ trợ đắc lực từ đôi tay nướng thịt điêu luyện của Botan và sức mạnh của Kaiho, một bữa tiệc nướng BBQ ngoài trời đã được thắp lửa. Bộ đôi nhí Naomi và Koneko được giao trọng trách làm ``phóng viên'', cần mẫn ghi chép lại từng khoảnh khắc giao lưu dở khóc dở cười giữa hai thế hệ: từ màn đọ lực nắm bóp lon nước của Kanata, vẻ ngượng ngùng đáng yêu của Towa trước Nene, đến trận đại chiến truyện cười thất bại giữa Polka và Ryuuhou. Khép lại trong tiếng cười giòn giã và bức ảnh chụp chung ấm áp, những tia sáng hy vọng tiếp tục âm thầm sưởi ấm căn phòng tầng tám trong màn đêm tĩnh lặng.
}

\pov{Hikawa Naomi}

Naomi có một cuốn sổ tay mới.

Đó là cuốn sổ bìa da màu nâu hạt dẻ nhỏ xíu, vừa vặn nằm gọn trong hai bàn tay của Naomi, do chị Suzuno tỉ mỉ dùng giấy tái chế đóng lại và tặng cho Naomi vào tối hôm qua. Đi kèm với cuốn sổ là một cây bút chì bảy màu có gắn cục gôm hình chú gấu trúc nhỏ ở đuôi.

Chị Suzuno bảo, hôm nay ký túc xá nhà mình sẽ có một sự kiện vô cùng náo nhiệt: một bữa tiệc nướng thịt ngoài trời quy tụ đông đủ nhất từ trước đến nay, có cả các chị bên Thế hệ thứ Tư lẫn Thế hệ thứ Năm cùng tới chơi. Vì thế, chị Suzuno giao cho Naomi cùng bạn Koneko một nhiệm vụ hết sức đặc biệt: làm ``tổ phóng viên thường trú'' của gia đình Hikawa, đi khắp nơi quan sát và ghi chép lại mọi chuyện xảy ra.

Koneko thích thú lắm. Cậu bạn nhỏ đeo chiếc túi đeo chéo hình mặt mèo, tay ôm theo chiếc kẹp bảng nhỏ để làm trợ lý đắc lực cho Naomi. Cả hai đứa chạy lăng xăng quanh nhà từ sáng sớm để bắt đầu công việc điều tra của mình.

Chiều thứ Bảy giữa tháng Bảy, bầu trời Kawachi trong xanh không một gợn mây. Nắng hè vàng rực rỡ như mật ong rót xuống từng ngọn cây phong ven đường, nhưng trên sân thượng tầng bảy của tòa ký túc xá, những cơn gió lộng từ dãy núi phía xa thổi về mát rượi, xua tan đi cái oi ả của buổi trưa.

Anh Kuon bảo, sân thượng này là nơi tuyệt vời nhất trên đời để ăn đồ nướng, vì khói than sẽ bay thẳng lên trời mà không làm ám mùi vào rèm cửa phòng ngủ của ai cả.

Và thế là, chiến dịch chuẩn bị cho đại tiệc BBQ chính thức bắt đầu.

\ct

\begin{center}
    \uline{Trích sổ tay phóng viên Naomi\\Trang 1: Phân chia lực lượng.}
\end{center}

Để chuẩn bị cho bữa tiệc đông người, anh Kuon đã chia mọi người ra làm ba nhóm khác nhau. Naomi và Koneko bí mật bám theo từng nhóm để ghi chép.

Nhóm thứ nhất là \textbf{Đội Nướng Thịt Sân Thượng}, gồm có anh Kuon, chị Botan và chị Kaiho.

Chị Botan trông rất tháo vát. Chị ấy mặc chiếc áo thun trắng đơn giản, bên ngoài đeo một chiếc tạp dề màu rằn ri, mái tóc xám buộc gọn gàng. Chị Botan dùng kẹp kim loại xếp từng thanh than hoa không khói vào chiếc lò nướng lớn bằng inox, vừa làm vừa cười vui vẻ:

``Kuon-senpai yên tâm đi! Khoản nhóm lửa với canh nhiệt độ thịt nướng cứ để em lo. Đảm bảo thịt bò wagyu hôm nay nướng lên vừa chín tới, mềm ngọt mà không bị khô đâu!''

Anh Kuon đứng bên cạnh, tay cầm quạt nan phẩy phẩy phụ họa, vẻ mặt nhẹ nhõm: ``Có Botan-chan ở đây anh yên tâm hẳn. Chứ để mấy đứa kia nhóm lửa chắc lát nữa cháy đen thui cả vỉ mất.''

Còn chị Kaiho thì phụ trách những việc nặng nhất. Chiếc vỉ nướng bằng gang đúc nguyên khối nặng chịch mà hai người lớn khiêng còn thấy vất vả, chị Kaiho chỉ dùng một tay nhấc bổng lên, đặt cái *KENG* vững như bàn thạch lên mặt bếp. Sau đó, chị lại vác thêm hai bao than củi ba mươi cân trên vai bước thoăn thoắt lên cầu thang mà trán không hề vương một giọt mồ hôi.

Chị Botan trố mắt nhìn chị Kaiho, hai mắt sáng rực vẻ bái phục: ``Kaiho-senpai khỏe thật đấy\dots\ Lực tay này chắc nhấc bổng cả cái xe máy lên cũng nhẹ tênh ấy nhỉ?''

Chị Kaiho chỉ khẽ nhếch mép, phủi bụi than trên tay: ``Xe máy thì không nói, nhưng nếu lát nữa đứa nào đùa giỡn làm đổ tàn than ra sàn gỗ thì chị cho đứa đó quét dọn sân thượng ba ngày. Thế thôi.''

Naomi nấp sau chậu cây xương rồng cùng Koneko, vội vã ghi vào sổ: *\uline{Chị Kaiho rất khỏe, nhưng cũng rất nghiêm khắc.}*

\ct

Nhóm thứ hai là \textbf{Đội Bếp Phụ Tầng Sáu}, đóng quân ở gian bếp lớn, gồm có chị Suzuno, chị Lamy và chị Watame.

Không khí ở đây dịu dàng và thơm tho hơn nhiều so với sân thượng đầy mùi khói than.

Chị Lamy đảm nhận khâu chuẩn bị đồ uống. Trên bàn đá cẩm thạch xếp ngay ngắn hàng chục chiếc ly thủy tinh cao cổ sáng loáng. Chị ấy cắt những lát chanh vàng, dâu tây và lá bạc hà thả vào bình thủy tinh lớn, rồi rót thêm trà đen ướp lạnh cùng nước khoáng có ga. Chị Lamy bảo đây là món trà hoa quả giải nhiệt mùa hè, người lớn hay trẻ con như Naomi và Koneko đều uống được. Chị còn đưa cho Naomi và Koneko mỗi đứa một quả dâu tây ngọt lịm nữa!

Trong khi đó, chị Watame thì đang tỉ mỉ xiên từng miếng ớt chuông, nấm hương và ngô ngọt vào que tre. Chị Watame đeo chiếc tạp dề in hình hoa cúc, vừa làm vừa quay sang cười với chị Suzuno:

``Suzuno-chan ơi, lát nữa nướng rau củ nhớ để riêng một góc vỉ sạch nhé! Đừng để dính mỡ thịt nhiều quá, ăn ngọt thanh tự nhiên mới ngon!''

Chị Suzuno đang dùng chiếc cân nhỏ để đong gia vị cho món nước sốt BBQ đặc biệt, nghe thấy thế bèn mỉm cười lễ phép: ``Chị Watame yên tâm đi ạ. Em đã làm riêng một âu sốt táo lê thanh ngọt chuyên dùng cho ngô nướng và nấm rồi, thơm lắm ạ!''

Naomi hí hoáy viết vào sổ: *\uline{Chị Watame rất thích ăn ngô nướng. Nước sốt của chị Suzuno thơm mùi táo lê.}*

\ct

Và cuối cùng là \textbf{Đội Phụ Việc Ngoài Sân} -- hay nói đúng hơn là những người bị chị Kaiho mời ra khỏi bếp vì quá ồn ào.

Nhóm này gồm có chị Luna, chị Polka và chị Ryuuhou.

Naomi và Koneko tận mắt chứng kiến cảnh lúc đầu giờ chiều: Chị Luna tò mò với lấy hũ mật ong trên kệ nhưng sơ ý làm trượt tay đổ ra bàn; chị Polka nổi hứng trổ tài tung hứng mấy quả cà chua bi rồi suýt làm rơi trúng âu nước sốt; còn chị Ryuuhou thì cứ loay hoay cầm chiếc kẹp gắp thịt gõ lách cách liên hồi làm ai cũng phân tâm.

Kết quả là chị Kaiho từ trên sân thượng bước xuống, khoanh tay đứng ở cửa bếp, giọng trầm xuống: ``Ba người ra ngoài sân phụ lau bàn ghế với xếp đĩa đi. Ở trong này chật chội, vướng chân người khác.''

Thế là ba người đành lủi thủi kéo nhau ra góc sân thượng, vừa lau bàn vừa đưa mắt nhìn về phía bếp than với vẻ thèm thuồng.

Chị Luna ngồi xuống ghế xếp, phụng phịu chu mỏ: ``Ư ư\~{}! Bất công ghê á-nanora! Em cũng muốn phụ mọi người nướng thịt mà-nanora!''

Chị Polka vừa lau bàn vừa thở dài: ``Định biểu diễn xiếc tạo không khí cho vui thôi mà, ai ngờ bị đuổi thẳng cánh cò bay thế này\dots''

Chị Ryuuhou cầm xấp đĩa giấy chia ra bàn, lắc đầu cười trừ: ``Thôi đi hai bà chị, cãi lời Kaiho-nee thì chỉ có nước nhịn ăn. Cứ ngoan ngoãn dọn bàn đi lát nữa tha hồ mà ăn.''

Koneko kéo kéo tay áo Naomi, chỉ vào chị Luna đang phồng má rồi thì thầm: ``Chị Luna trông buồn cười ghê á.''

Naomi gật gù, nắn nót ghi lại: *\uline{Ba chị bị phạt vì quá lộn xộn. Bạn Koneko bảo chị Luna giống chim cánh cụt giận dỗi.}*

\ct

Khoảng năm giờ chiều, khi những dải mây màu cam đỏ bắt đầu trải dài nơi chân trời phía tây, bữa tiệc nướng chính thức khai mạc.

Mùi thịt nướng xèo xèo bốc lên ngào ngạt, hòa quyện với mùi thơm của bơ tỏi, hành tây nướng và nước sốt hoa quả làm bụng Naomi và Koneko réo lên ùng ục. Tất cả mọi người cùng kéo nhau lên sân thượng, vây quanh chiếc bàn dài bằng gỗ sồi đã được bày biện la liệt những đĩa thức ăn hấp dẫn.

Và đây cũng là lúc những màn giao lưu vui nhộn giữa hai thế hệ bắt đầu diễn ra.

\stpt{-Sự kiện thứ nhất-}
\stcap{Trận đọ lực bóp lon nước ngọt}

Chuyện bắt đầu khi chị Botan đang ngồi uống nước ngọt có ga ở góc bàn, đưa mắt nhìn sang chị Kanata đang bưng một đĩa thịt bò lớn bằng một tay rất nhẹ nhàng.

Chị Botan cười tươi, mở lời trêu: ``Này Kanata-san! Lực tay của chị nổi tiếng khắp công ty từ lâu rồi đấy nhé. Hay là nhân dịp đông đủ thế này, hai đứa mình thử đọ xem tay ai khỏe hơn chút không?''

Chị Kanata đang gắp thịt, nghe thấy thế bèn quay sang, cười toe toét: ``Đọ kiểu gì thế Botan-chan? Chẳng lẽ lại vật tay ngay trên bàn ăn à?''

Chị Botan giơ chiếc lon nước ngọt bằng nhôm rỗng lên: ``Bóp lon nước rỗng thôi! Xem ai bóp gọn gàng hơn!''

Cả bàn tiệc lập tức hào hứng hùa theo. Chị Polka và chị Nene đứng bật dậy vỗ tay cổ vũ: ``Hay đấy! Thi đi hai bà ơi!''

Chị Botan thử trước. Bàn tay khỏe khoắn của cô nàng siết mạnh lấy chiếc lon nhôm. *RẮC!* Chiếc lon kim loại lập tức bẹp dí, biến dạng thành một khối tròn méo mó gọn lỏn trong lòng bàn tay. Chị Botan đặt khối nhôm xuống bàn, cười đắc ý: ``Thế nào?''

Mọi người vỗ tay khen ngợi.

Chị Kanata chỉ mỉm cười nhẹ nhàng. Chị ấy cầm lấy một chiếc lon rỗng khác, đặt vào giữa lòng bàn tay phải của mình.

Chị không hề gồng mình, cũng chẳng lấy đà gì cả. Chỉ thấy những ngón tay của chị khẽ siết lại một cái thật tự nhiên.

*RẮC\dots\ XẸP!*

Một tiếng động khô khốc vang lên. Khi chị Kanata mở bàn tay ra, chiếc lon nhôm không chỉ bị bẹp, mà nó đã bị ép phẳng lì hoàn toàn, mỏng dính chẳng khác nào một tấm kim loại lót ly!

Toàn bộ sân thượng im lặng mất mấy giây.

Chị Botan tròn xoe mắt nhìn chiếc lon bị dát mỏng, rồi nhìn sang bàn tay nhỏ nhắn của chị Kanata, sau đó ôm bụng cười ngả nghiêng: ``Trời đất ơi! Quả nhiên lời đồn không sai chút nào! Lực nắm thế này thì ai mà đọ lại được chứ! Em xin chịu thua sát nút luôn!''

Chị Kanata ngượng ngùng cười xòa: ``Có gì đâu mà, tại cái vỏ lon này nó mỏng sẵn đấy chứ\dots''

Naomi và Koneko ghé đầu vào nhau quan sát, nhanh tay vẽ một hình minh họa chiếc lon dẹp lép vào sổ kèm dòng chú thích: *\uline{Tay chị Kanata khỏe như kìm ép sắt. Bạn Koneko bảo không nên chọc chị ấy giận.}*

\ct

\stpt{-Sự kiện thứ hai-}
\stcap{Lời khen chân thành và đôi má ửng hồng}

Ở phía bên kia bàn ăn, một tình huống đáng yêu khác lại đang diễn ra giữa chị Towa và chị Nene.

Chị Towa đang cẩn thận dùng kẹp gắp từng miếng xúc xích và thịt nướng chín vàng chia đều vào đĩa cho mọi người xung quanh, miệng vẫn luôn nhắc: ``Mọi người ăn từ từ thôi kẻo bỏng đấy nhé.''

Chị Nene lon ton bưng đĩa lại gần, cười tít mắt: ``Cảm ơn Towa-senpai nha! Senpai chu đáo ghê luôn á!''

Chị Towa hơi khựng lại, quay mặt đi chỗ khác, cố giữ giọng bình thản: ``Có gì đâu, tiện tay gắp thôi mà. Mọi người cứ mải nói chuyện để thịt cháy xém hết cả ra kia kìa, ăn nhanh lên giùm cái!''

Chị Nene thấy điệu bộ ngượng ngùng ấy thì thích thú lắm, chồm người lại gần hơn, hai mắt long lanh: ``Hi hi, senpai cứ làm bộ hoài à! Rõ ràng là senpai thấy đĩa của em hết sạch nên mới gắp cho em miếng to nhất đúng không? Lại còn chu đáo thổi bớt khói nữa chứ! Towa-senpai lúc nào cũng tốt bụng và dễ thương hết trơn!''

Gương mặt chị Towa lập tức đỏ bừng lên tận mang tai. Chị ấy cuống quýt xua tay, lắp bắp: ``N-Này nhé Nene! Đừng có nói linh tinh thế chứ! Chị\dots\ chị chỉ làm tròn bổn phận thôi mà! Mau ăn đi, bớt nịnh nọt giùm tôi cái!''

``Không nịnh đâu nè, em nói thật lòng mà!'' Chị Nene cười giòn tan, tiện tay xiên một miếng đào ngọt đút cho chị Towa.

Bị miếng đào ngọt lịm chặn lời, chị Towa vừa nhai vừa lúng túng nhìn đi nơi khác, nhưng khóe môi thì không giấu nổi một nụ cười vui vẻ.

Naomi gật gù ghi lại: *\uline{Chị Towa rất hay ngại ngùng khi được khen, nhưng chị ấy luôn là người quan tâm gắp đồ ăn cho mọi người trước tiên.}*

\ct

\stpt{-Sự kiện thứ ba-}
\stcap{Trận chiến truyện tiếu lâm bất thành}

Trong khi đó, ở khu vực bàn mây gần lan can, chị Polka và chị Ryuuhou lại đang tranh luận xem ai có thể kể chuyện cười vui hơn.

Chị Polka đứng lên cạnh bàn, hắng giọng làm vẻ trịnh trọng: ``Này mọi người, tôi có câu đố này hay lắm nè! Đố biết, con vịt đi bằng hai chân thì gọi là con gì?''

Chị Watame ngây thơ chớp mắt: ``Là con vịt bình thường đúng không bà?''

``Sai rồi! Đó là vịt Donald!'' Chị Polka cười sảng khoái, rồi đố tiếp: ``Thế con chuột đi bằng hai chân thì gọi là gì?''

Chị Kson nhấp ngụm nước ngọt, đáp bừa: ``Chuột Mickey?''

``Chuẩn luôn!'' Chị Polka đập tay xuống bàn: ``Thế còn con vịt đi bằng hai chân thì gọi là gì nào?!''

Chị Nene nhanh nhảu: ``Vịt Donald tiếp đúng không bà?''

``Sai bét tè lè nhè! Con vịt nào mà chẳng đi bằng hai chân hả mấy bà! Hahaha!''

Chị Polka tự cười nghiêng ngả một mình, nhưng xung quanh\dots\ cả bàn tiệc chỉ chớp chớp mắt nhìn nhau. Tiếng gió hè thổi qua mang theo tiếng lá cây xào xạc. Không một ai cười nổi. Chị Kson chỉ nhếch mép thở dài, còn chị Kaiho thì lắc đầu ngao ngán.

Chị Ryuuhou thấy vậy bèn cười khẩy: ``Trời ơi chị Polka ơi, đố mẹo con nít thế ai mà cười cho nổi! Để em kể chuyện này cười rách mép luôn nè!''

Chị Ryuuhou làm bộ mặt nghiêm túc: ``Có một cái que tăm đang đi dạo trên đường thì thấy một con nhím đi ngang qua. Cái que tăm mới giật mình bảo: `Ủa, xe buýt chạy tuyến này từ bao giờ thế nhỉ?'\dots\ Hết rồi đó! Mắc cười không?''

Một giây trôi qua. Hai giây trôi qua.

Vẫn không một tiếng cười nào vang lên. Chỉ có tiếng than củi nổ *lách tách* trên lò nướng.

Chị Ryuuhou đứng hình, nụ cười trên môi méo xệch đi.

Đúng lúc đó, anh Kuon bưng đĩa thịt đi ngang qua, buông một câu nhận xét tỉnh bơ: ``Hai đứa bớt kể chuyện cười lại giùm anh cái. Nhạt đến mức gió hè cũng muốn đóng băng luôn rồi đấy.''

*PHỤT!*

Chị Kson và chị Botan lập tức ôm bụng cười sặc sụa, tiếng cười như xé tan không khí ngượng ngùng vừa rồi. Cả bàn tiệc bỗng chốc vỡ òa trong những tràng cười giòn giã—không phải vì hai câu chuyện đùa dở tệ kia, mà là vì vẻ mặt bất lực của anh Kuon!

Naomi khúc khích cười, nắn nót viết: *\uline{Chị Polka và chị Ryuuhou kể chuyện không ai cười, nhưng anh Kuon trêu một câu là cả nhà cười rần rần.}*

\ct

\stpt{-Sự kiện thứ tư-}
\stcap{Góc xích đu và đĩa thịt nướng yêu thương}

Giữa những tiếng cười đùa ồn ã ấy, Naomi và Koneko để ý thấy ở góc khuất gần giàn hoa giấy, chị Delutaya đang ngồi một mình trên chiếc xích đu gỗ.

Chị Delu mặc chiếc áo hoodie màu tím than quen thuộc, hai tay ôm chiếc đĩa nhỏ đựng hai miếng thịt nướng và nửa bắp ngô ngọt, đôi mắt tròn xoe lặng lẽ dõi theo mọi người. Dù nụ cười vẫn nở trên môi, nhưng đôi vai hơi co lại cho thấy chị ấy vẫn còn chút ngập ngừng trước một bữa tiệc đông đúc náo nhiệt thế này.

Chị Lamy từ bàn tiệc nhẹ nhàng bước lại gần, mang theo hai ly nước ép táo mát lạnh. Chị ngồi xuống khoảng trống bên cạnh chiếc xích đu, đưa ly nước cho chị Delu, cất giọng dịu dàng: ``Thịt nướng có vừa miệng bà không, Delu-chan?''

Chị tóc xanh lá khẽ giật mình, rồi nở nụ cười ngượng ngùng: ``Ừm\dots\ ngon lắm bà. Lâu lắm rồi tôi mới được ăn một bữa tiệc đông vui thế này. Chỉ là\dots\ thấy đông người quá, tự nhiên tôi hơi bị ngợp một chút thôi.''

Chị Lamy mỉm cười thấu hiểu, đưa tay vỗ nhẹ lưng bạn mình: ``Không sao đâu, cứ thoải mái đi. Bà nhìn xem kìa, ở đây toàn là người một nhà với bạn bè thân thiết cả, có ai xa lạ đâu nào.''

Đúng lúc đó, chị Watame và chị Kson từ bàn nướng cũng mang đĩa cánh gà nướng sốt mật ong thơm lừng chạy lại: ``Này Delu-chan! Ăn thử cánh gà nướng mật ong đi nè! Wata-chan vừa nướng xong còn nóng hổi luôn á!'' chị Watame cười tươi rói.

Chị Kson cũng quàng tay qua vai chị Delu, vỗ bồm bộp: ``Đúng đó! Ra bàn ngồi chung với tụi này cho xôm tụ, ngồi một góc gió thổi lạnh đấy!''

Koneko và Naomi cũng chạy lon ton lại, Koneko đưa cho chị Delu một xiên nấm nướng thơm phức: ``Chị Delu ăn nấm nướng này đi ạ, ngon lắm đó!''

Nhìn thấy ánh mắt ấm áp không một chút khoảng cách của tất cả mọi người, sự rụt rè trong lòng chị Delu dường như tan biến sạch sẽ. Chị ấy cắn một miếng nấm nướng, rồi bật cười thành tiếng bằng một nụ cười rạng rỡ, tự nhiên và hạnh phúc vô cùng.

Naomi mỉm cười, viết vào trang sổ: *\uline{Chị Delu đã cười rất tươi. Bạn Koneko bảo nấm nướng mang lại phép màu.}*

\ct

Trời về khuya, những ngọn đèn lồng giấy treo quanh sân thượng được thắp sáng lung linh, tỏa ra ánh sáng vàng ấm áp dưới bầu trời đêm đầy sao.

Sau khi mọi người đã ăn uống no nê và cùng nhau dọn dẹp sơ bộ, chị Suzuno khẽ mỉm cười quay sang Naomi:
``Naomi-chan, bản tin phóng viên của em đã hoàn thành chưa? Đọc cho cả nhà nghe một chút được không?''

Tất cả ánh mắt lập tức đổ dồn về phía hai cô bé.

Naomi hơi ngượng, nhưng được bạn Koneko đứng bên cạnh nắm tay cổ vũ, Naomi hít một hơi thật sâu, mở cuốn sổ tay ra đọc to từng dòng chữ: ``Bản tin đặc biệt của ký túc xá Hikawa, phóng viên Hikawa Naomi và trợ lý Koneko xin được báo cáo:

Hôm nay là một ngày tuyệt vời nhất mùa hè.

Kaiho-oneechan là người khỏe nhất vì có thể nhấc bổng vỉ nướng bằng một tay, nhưng chị không cho ai làm bẩn sàn nhà.

Botan-oneechan nướng thịt rất ngon và cười rất giòn tan.

Kanata-oneechan có thể nghiền nát vỏ lon nhôm phẳng lì như tờ giấy.

Towa-oneechan rất hay xấu hổ khi được khen, nhưng chị luôn chu đáo gắp thức ăn ngon cho mọi người.

Polka-oneechan và Ryuuhou-oneechan kể chuyện đố vui không ai cười, nhưng hai chị ăn thịt nướng rất nhanh.

Delutây-oneechan ban đầu ngồi một mình, nhưng sau đó đã cười rất nhiều vì có các chị và bạn Koneko chia đồ ăn.

Và cuối cùng\dots\ Onii-chan là người ăn nhiều thịt vụn nhất bữa tiệc, vì anh ấy bảo anh ấy phải `thử độc' trước khi mang ra cho cả nhà ăn!''

Naomi vừa dứt lời, toàn bộ sân thượng bùng nổ trong những tràng pháo tay và tiếng cười ròn rã.

Chị Botan đập bàn cười ngặt nghẽo, chị Towa thì ôm mặt che nụ cười ngượng nghịu, còn anh Kuon thì gãi đầu cười bất lực: ``Trời đất ơi hai đứa nhóc này\dots\ Ai dạy em mấy cái suy luận đó thế hả Naomi?''

Chị Kson đứng bật dậy, trên tay cầm chiếc máy ảnh polaroid chụp lấy liền:

``Hay lắm nhóc Naomi và Koneko! Báo cáo quá chuẩn xác không trượt phát nào! Nào mọi người ơi, đứng sát lại đây! Ngày vui thế này mà không làm một kiểu ảnh kỷ niệm thì phí cả đời!''

Thế là tất cả mọi người cùng ùa lại gần nhau bên bếp than hồng.

Anh Kuon đứng ở vị trí trung tâm, chị Kaiho khoanh tay đứng cạnh với nụ cười hiếm hoi nở trên môi. Các cô gái Gen 4 và Gen 5 quây quần xung quanh, người giơ tay chữ V, người ôm vai nhau cười toe toét. Naomi và Koneko được chị Suzuno cho ngồi ở hàng ghế đầu, hai đứa nhỏ cùng giơ cao cuốn sổ phóng viên đầy tự hào.

*TÁCH!*

Ánh đèn flash bừng sáng xé tan màn đêm.

Chiếc máy ảnh nhè nhẹ nhả ra một tấm phim nhỏ. Dưới ánh trăng hè vằng vặc, những nụ cười rạng rỡ, những ánh mắt lấp lánh và hơi ấm của tình bạn chân thành đã vĩnh viễn được khắc ghi trên mặt giấy bóng loáng.

Một bức ảnh lịch sử đầu tiên quy tụ trọn vẹn cả Thế hệ thứ Tư, Thế hệ thứ Năm và đại gia đình Hikawa.

\ct

Mười một giờ đêm.

Bữa tiệc đã tàn. Những vị khách quý của hai thế hệ đã lên chuyến tàu điện muộn trở về Tokyo, mang theo những nụ cười và kỷ niệm ngọt ngào.

Các thành viên trong nhà cũng đã mệt lử sau một ngày dài náo nhiệt, lần lượt chìm vào giấc ngủ say nồng trong phòng riêng. Cả tòa nhà ký túc xá trở về với vẻ tĩnh mịch quen thuộc.

Trước khi đi ngủ, chị Suzuno đã cùng bạn Koneko cẩn thận chuẩn bị một chiếc khay sứ nhỏ: bên trên là một đĩa thịt bò nướng thơm phức được giữ ấm cẩn thận, vài lát ngô ngọt và rau củ thanh đạm, cùng một ly trà hoa cúc dịu mát.

Koneko nhẹ nhàng bưng chiếc khay lên tầng tám, đặt ngay ngắn trước cửa căn phòng số 801, rồi gõ nhẹ ba tiếng *cốc, cốc, cốc* quen thuộc: ``Onee-chan\dots\ Mọi người để phần đồ nướng ngon lắm cho chị nè. Chị nhớ ăn cho ấm bụng nhé.''

Nói xong, cô bé khẽ mỉm cười rồi quay về phòng ngủ.

Một lúc sau, anh Kuon chậm rãi bước lên kiểm tra hành lang tầng tám theo thói quen mỗi đêm.

Khi bước chân anh dừng lại trước căn phòng, cánh cửa gỗ vốn luôn đóng kín lúc này đang hé mở một khe nhỏ. Ánh sáng vàng dịu êm le lói hắt ra sàn gỗ bóng loáng.

Và trên chiếc khay gỗ trước cửa phòng\dots

Chiếc đĩa sứ trắng tinh lúc này đã hoàn toàn trống trơn.

Từng miếng thịt nướng, từng lát ngô ngọt đều đã được ăn sạch sẽ không còn sót lại một vụn nhỏ. Chiếc đĩa đã được lau sạch bóng, và đặt ngay ngắn bên cạnh là một mẩu giấy nhớ màu hồng nhỏ xíu, được gấp thành hình một bông hoa tuyết vụng về nhưng vô cùng nắn nót.

Anh Kuon đứng lặng người trong bóng tối của hành lang.

Một luồng cảm xúc ấm áp dâng trào từ tận đáy lòng, lan tỏa khắp lồng ngực anh. Sống mũi anh cay cay, khóe môi khẽ cong lên một nụ cười nhẹ nhõm và bình yên nhất trần đời.

Em ấy đã ăn hết.

Em ấy đã đón nhận hơi ấm của bữa tiệc, đón nhận tấm lòng của mọi người, dẫu cho vẫn còn đang ẩn mình sau cánh cửa gỗ khép hờ.

Cơn bão lớn ngoài kia có thể vẫn còn nhiều thử thách, nhưng ở nơi này, tình yêu thương và sự kiên nhẫn chân thành đang từng ngày làm tan chảy đi lớp băng giá lạnh lẽo nhất.

Anh Kuon khẽ cúi người, nhẹ nhàng nhấc chiếc khay sứ lên tay, rồi thì thầm một câu thật khẽ vào màn đêm tĩnh mịch:

``Ngủ ngon nhé, Mike-chan.''

Phía sau khe cửa hẹp, ánh đèn vàng khẽ chớp tắt một nhịp, tựa như một lời chúc ngủ ngon không thành tiếng gửi lại cho người anh cả của ngôi nhà.

%==========================================TRUYỆN PHỤ=========================================
\omk

\pov{-}

Sáng hôm sau, sáu giờ ba mươi phút tại phòng khách tầng năm.

Kuon vừa pha xong một cốc cà phê đen nóng, ngồi xuống ghế sofa xoa xoa hai bên thái dương còn hơi mỏi sau một buổi tối dọn dẹp vất vả. Chiếc đồng hồ thông minh trên cổ tay anh bỗng chốc phát ra một tiếng *bíp* đanh gọn, ánh sáng xanh lam từ mặt kính quét lên một màn hình hologram ba chiều lơ lửng giữa không trung.

Game xuất hiện với biểu tượng khuôn mặt pixel nhấp nháy, cất giọng mỉa mai sắc lẹm nhưng đầy vẻ chuyên nghiệp: ``Chào buổi sáng, Kuon-user. Căn cứ theo yêu cầu giám sát an ninh và phân tích dữ liệu hành vi môi trường, tôi đã hoàn tất bản báo cáo tổng hợp âm học và sinh trắc học của sự kiện BBQ tối ngày hôm qua. Đề xuất anh nên xem qua trước khi đưa ra quyết định tổ chức thêm bất kỳ hoạt động tập thể tương tự nào trong tương lai.''

Kuon nhấp một ngụm cà phê, nhướng mày: ``Báo cáo gì thế? Lại số liệu tào lao của ông à?''

``Hoàn toàn là dữ liệu khoa học chính xác 100\%,'' Game lạnh lùng hiển thị hàng loạt biểu đồ cột và đồ thị hình sin chạy ngoằn ngoèo:

\txtbx{
    \textbf{BÁO CÁO PHÂN TÍCH DỮ LIỆU SỰ KIỆN BBQ (16/07/2022)} \\[4pt]
    - \textbf{Tổng số lần Omaru Polka phát âm thanh vượt ngưỡng 85 decibel:} 47 lần (tương đương tiếng ồn của máy cắt cỏ hoạt động liên tục). \\
    - \textbf{Tổng số lần Himemori Luna sử dụng hậu tố ngữ pháp `nanora':} 83 lần (tần suất trung bình: 1,38 nanora/phút). \\
    - \textbf{Tổng số lần Hikawa Kaiho nhíu mày và phát ra tín hiệu đe dọa kỷ luật:} 126 lần (nguy cơ kích hoạt nắm đấm thép đạt mức báo động đỏ 4 lần). \\
    - \textbf{Độ giãn nở dạ dày của Tsunomaki Watame:} Vượt ngưỡng an toàn sinh học của loài cừu 142\% do tiêu thụ quá mức ngô nướng. \\
    - \textbf{Lượng thịt nướng do Kuon-user tiêu thụ dưới danh nghĩa `thử độc':} 650 gam (vượt 200\% chỉ tiêu ăn kiêng chống béo bụng).
}

Kuon sặc cà phê, ho sù sụ: ``Này! Cái dòng cuối cùng là do ông tự bịa ra đấy à?! Tôi nếm thử đồ ăn trước xem có chín đều không thôi mà!''

Game nhấp nháy biểu tượng mặt cười khinh bỉ: ``Dữ liệu từ cảm biến nhịp tim và enzym tiêu hóa không biết nói dối, Kuon-user à. Tuy nhiên, một chỉ số bất thường đã được ghi nhận: Mặc dù mức độ hỗn loạn âm thanh tương đương một vụ nổ cấp ba, nhưng nồng độ hormone cortisol gây căng thẳng trong máu của cậu lại giảm 18\% so với ngày thường, và nhịp tim trung bình duy trì ở mức thư thái nhất trong vòng sáu tháng qua.''

Game dừng lại một nhịp, rồi buông một câu kết luận hiếm hoi không mang tính châm chọc: ``Đánh giá chung của hệ thống: Mức độ thành công của sự kiện đạt 98,5\%. Gia đình Hikawa và các đối tượng liên quan đã sống sót qua thảm họa một cách kỳ diệu. Hết báo cáo.''

Kuon nhìn màn hình hologram dần tắt lịm, rồi bật cười thành tiếng. Anh tựa lưng vào thành ghế êm ái, ngắm nhìn ánh nắng ban mai rực rỡ đang chiếu qua khung cửa sổ phòng khách.

Một ngày mới lại bắt đầu, bình yên và tràn ngập niềm vui.

\end{document}
